% ===================== 第 10 章 =====================
\section{ForesightBlindspot 评测体系}

\subsection{为什么需要时间切片基准}

盲点发现的价值在于\textbf{提前看见后来才被证实的缺口}，普通检索质量基准度量不了这种能力。ForesightBlindspot 因此采用时间切片设计：每个案例是一个争议问题加一个截止日，系统只能读取截止日前的封闭语料；金标准盲点来自「后来被高质量研究证实、截止日前已有合理先兆、且不依赖不可预见事件」的缺口。协议内置泄漏控制：封闭语料、禁止未来引用、变量可匿名化，并明确披露基础模型参数可能包含截止日后知识这一无法完全消除的风险。

\subsection{五条基线、六项消融：每一级只加一个能力}

五条基线全部由同一套生产组件接线而成，而非另写「对照组实现」，因此相邻两级的差异恰好是一个机制：

\begin{enumerate}[leftmargin=1.6em, itemsep=2pt]
\item Single-Agent Deep Research：单席、无协议、无记忆、无证据门；
\item Fixed Multi-Agent Debate：七席自由辩论，无每席记忆、无证据门；
\item Council + Linear Context：七席共享一份无差别转录；
\item Council + MemoBrain without Evidence Engine：有每席记忆但移除证据门；
\item Full Poliscope：完整系统。
\end{enumerate}

六项消融各移除一个机制：独立预承诺、对抗性证伪者、证据审计员、辩证折叠、证据谱系、MemoBrain。指标分五族：盲点（召回/精度/高影响召回/可行动性/预见距离）、证据（引用存在、引用蕴含、独立性、因果过度推断）、记忆（压缩比、骨架保留、异议保留、长程漂移）、协作（角色多样性、信念更新、虚假共识率）与成本（每个有效盲点成本、每条入图发现 token）。

\subsection{评测管道与生产系统同构}

\code{packages/evaluation/harness.py} 从生产的编排器、交付器与记忆组件直接接线出全部基线与消融，打分函数复用生产侧的因果越级策略、引用蕴含验证器与谱系聚类，避免「评测里一套、生产里一套」。管道第一次运行就抓出两个基线构造错误，并作为方法教训写进论文：单 Agent 基线曾按枚举序误选证伪者席位（恰好在脚本素材上答案最强），导致单 Agent 得分与完整议会持平；最终复判曾无条件遍历七席、给从未参与的席位铸造占位判断。两者修复后，回归测试断言「单席覆盖严格低于多席」，防止基线阶梯再次坍塌。

\subsection{真实模型实验：结果连同不利证据一起报告}

真实模型对照使用 DeepSeek-V4-Flash 官方兼容端点：首轮 11 变体扫描遭遇端点限流，12 个被污染的 run 直接弃用；随后对五个关键变体各运行 3 次取均值。盲点裁判同步升级为语义评判（\code{semantic\_blindspot.py}，复用同一模型网关）——关键词子串匹配在自由文本下会把「效应方向可能反向」误判为零命中，这是测量仪器失效，不是系统能力为零。表~\ref{tab:live} 为三轮均值。

\begin{table}[ht]
\centering\small
\caption{真实模型五变体三轮均值（语义裁判；脚本 \code{scripts/live\_ablation.py}，原始数据 \code{docs/paper/data/}）}
\label{tab:live}
\begin{tabularx}{\linewidth}{@{}L{3.9cm} c c c c c@{}}
\toprule
\rowcolor{accent}
\tblhead 变体 & \tblhead 语义召回 & \tblhead 语义精度 & \tblhead F1 & \tblhead 空槽 & \tblhead 事件数 \\
\midrule
完整 Poliscope                       & 0.286 & \textbf{1.000} & 0.444 & 12.3 & 118.3 \\
单 Agent 基线                        & 0.809 & 0.544 & 0.651 & 1.0  & 55.7 \\
去独立预承诺                         & \textbf{1.000} & 0.490 & 0.658 & 1.0  & 160.0 \\
去对抗性证伪者                       & 1.000 & 0.406 & 0.577 & 1.7  & 160.3 \\
去证据审计员                         & 0.952 & 0.586 & 0.726 & 1.3  & 148.0 \\
\bottomrule
\end{tabularx}
\end{table}

\begin{findingbox}
\boxtag{\color{okgreen}三个可以站住的发现}
\textbf{1. 语义裁判修复了测量。}关键词版对完整系统报零召回，语义版给出可读、非退化的分布，精度形成真实梯度。\\
\textbf{2. 这是稳定的精度—召回权衡，不是噪声。}完整系统精度满分、召回最低（三次都恰好覆盖 7 个金概念中的 2 个）；所有简化变体召回 0.81–1.00、精度降到 0.41–0.59；去预承诺的满召回三次复现一致。差距由空槽数解释：完整系统平均 12.3 个结构化提名产不出来，简化变体仅 1.0–1.7 个——议会治理机制以召回换精度，而 flash 级模型扛不住完整协议负担时，召回代价主导了 F1。\\
\textbf{3. 数据层确定性评分与模型无关。}引用蕴含 0.5、证据独立性 0.667、异议保留 1.0 在全部 15 个 run 中完全一致。
\end{findingbox}

\subsection{诚实边界}

时间切片封闭语料尚未采集，独立的 B\_test 测试集暂缺；双人标注加仲裁的一致性骨架（Cohen's $\kappa$ / Krippendorff's $\alpha$，\code{agreement.py}）已实现，但尚无标注数据。脚本化确定性测量（B\_dev，无模型、无网络）只能隔离各机制的结构贡献，协议与记忆对盲点\emph{整合}的收益必须在足以支撑完整协议的模型上重新观测。这些缺口在论文与本白皮书中同等对待：标为未来工作，不用脚本化数字冒充真实模型结论，也不把裁判不可达时的缺失值静默写成零。
